\chapter*{Tóm tắt}
\markboth{Tóm tắt}{Tóm tắt}
\addcontentsline{toc}{chapter}{Tóm tắt}
\onehalfspacing
\vspace{1.0cm}

Truy xuất ảnh có điều kiện (\emph{Composed Image Retrieval} -- CIR) tìm ảnh đích từ một truy vấn gồm ảnh tham khảo và văn bản mô tả chỉnh sửa. Mọi phương pháp CIR trước đây đều giả định truy vấn chỉ có một ảnh, trong khi sản phẩm thời trang trên các trang thương mại điện tử được trình bày từ nhiều góc nhìn. Bài báo \textbf{FashionMV} (Yuan và cộng sự, 2026) gọi sự lệch pha này là \emph{View Incompleteness}, chính thức định nghĩa bài toán \textbf{Multi-View CIR} ở mức sản phẩm, xây dựng bộ dữ liệu FashionMV (127 nghìn sản phẩm, 472 nghìn ảnh, hơn 220 nghìn triplet) và đề xuất mô hình \textbf{ProCIR} trên nền MLLM Qwen3.5-0.8B với ba cơ chế: hội thoại hai lượt, căn chỉnh dựa trên chú thích và chuỗi suy luận.

Báo cáo Thực tập 1 tổng hợp có hệ thống bài báo này cùng cơ sở lý thuyết liên quan, sau đó tái lập thực nghiệm bằng checkpoint và mã đánh giá công khai của tác giả. Do nguồn ảnh gốc không khả dụng, học viên xây dựng quy trình thay thế bằng các bản sao trên HuggingFace và chạy đánh giá trên GPU miễn phí của Kaggle. ProCIR tái lập đạt R@5/R@10 = 74,42/85,25 trên toàn bộ 5.188 triplet DeepFashion (bài báo: 89,2/94,9) và 87,93/94,53 trên mẫu 1.500 triplet Fashion200K (bài báo: 77,6/86,6). Hai sai lệch ngược chiều được truy nguyên đến nguyên nhân định lượng được: ảnh DeepFashion thay thế chỉ có độ phân giải $224 \times 224$ (ít token thị giác hơn khoảng 5 lần), và gallery Fashion200K bị thu nhỏ khoảng 4,5 lần khi đánh giá trên mẫu con. Hai baseline zero-shot CLIP và FashionCLIP với hợp nhất muộn được đánh giá trên đúng tập triplet và gallery của ProCIR; thứ tự CLIP $<$ FashionCLIP $<$ ProCIR được giữ nhất quán, với khoảng cách hơn 19 điểm R@5 giữa ProCIR và cấu hình baseline tốt nhất trên DeepFashion. Báo cáo cũng chỉ ra rằng giao thức đánh giá gốc giữ sản phẩm nguồn trong gallery: với baseline, sản phẩm nguồn chiếm hạng 1 ở gần như mọi truy vấn, làm R@1 giảm gần về 0 -- một chi tiết cần được chuẩn hoá khi so sánh các phương pháp.

Trong một thực nghiệm bổ sung sau khi có đủ ảnh của cả ba bộ dữ liệu, học viên cài đặt lại phần huấn luyện và huấn luyện ProCIR trên 25\% dữ liệu với LoRA. Trên toàn bộ tập validation, cấu hình MT+Align đạt trung bình R@5 = 72,47 so với 71,26 của single-turn, lặp lại xu hướng ablation của bài báo. Trên cùng bộ ảnh, checkpoint của tác giả đạt R@5 = 89,26 trên DeepFashion và 75,82 trên FashionGen, khớp với số công bố, nhưng chỉ đạt 67,71 trên Fashion200K (bài báo: 77,6), cho thấy ảnh Fashion200K đang dùng không tương đương với ảnh của tác giả.

\bigskip
\noindent\textbf{Từ khóa:} truy xuất ảnh có điều kiện, Multi-View CIR, mô hình ngôn ngữ lớn đa phương thức, truy xuất ảnh thời trang, tái lập thực nghiệm.
